% =============================================================================
% 049 / 068
% Status: raw
% =============================================================================
% Notes:
%
% Source question:
% 像越南或韓國那樣採用全拼音文字，有哪些缺點？
% 為什麼其他語言完全採用拼音是沒有問題的，(前)漢字圈諸語言卻不行？
% =============================================================================

\chapter[像越南或韓國那樣採用全拼音文字，有哪些缺點？ 為什麼其他語言完全採用拼音是沒有問題的，(前)漢字圈諸語言卻不行？]{像越南或韓國那樣採用全拼音文字，有哪些缺點？ \\[0.35em] 為什麼其他語言完全採用拼音是沒有問題的，(前)漢字圈諸語言卻不行？}
\qaanswer{
漢字圈文化的語言都受到中古漢語、近古漢語強烈的影響與滲透。漢字詞彙大量地置換了這些地區的原生詞彙，漢字已經作為詞根的形式深度紮根與這些語言的詞彙中。而漢字詞根比較的簡練，導致重音律高，全拼音化的結果會造成同音異議詞變多，尤其是朝鮮語，拋開漢字後閱讀都會造成一定的障礙。再加上東亞語言詞彙的虛詞詞根部分特徵不明顯，全漢字導致虛實詞難以區分，全拼同樣會造成虛實詞難以區分。這一點採取實詞用漢字，虛詞用假名的日語構建是最好的。除此之外，漢字圈文化中詞性標記整理的相對好的是日語和朝鮮語，前者雖然有大量的訓讀但問題是音素太少，後者則是全文讀系統，這都導致全拼情況下記憶詞彙會變的困難。但這些其實都不是最嚴重的問題，真正的問題在於全拼後導致文化的割裂，今人無法很好的繼承祖先的文明傳承。大量文明資訊因不懂漢字而造成根源上的丟失。這一點朝鮮半島人和越南人深有體會，這兩地文字改革後一直沒有出現著名的文學作品也從側面反映了這一點。
}

